\section{Testing signs in maximal tensor products}
\label{sec:tensor}
The goal is to prove Theorem~\ref{thm:tensor-intro}: the two signs of
one selfadjoint element will survive in every irreducible representation.
Such representations of a maximal tensor product need not be spatial.
The following factor argument and the pure-state excision method of
Akemann--Anderson--Pedersen \cite[Proposition~2.2]{AAP} let us transfer
the real evaluations of Lemma~\ref{lem:gaussian} without a spatial
representation hypothesis. We include the required excision proof.

\begin{lemma}[Products in commuting factors]\label{lem:factors}
Let \(M_1,\ldots,M_n\subseteq\mathcal B(H)\) be mutually commuting unital von Neumann factors. If \(e_j\in M_j\) is positive and \(\|e_j\|=1\), then
\[
 \left\|\prod_{j=1}^ne_j\right\|=1.
\]
\end{lemma}
\begin{proof}
First, a product of nonzero projections \(p_j\in M_j\) is nonzero. Induct on \(n\). By induction \(q=p_2\cdots p_n\) is a nonzero projection commuting with \(M_1\). The restriction \(M_1\to\mathcal B(qH)\) is a nonzero normal representation. Its kernel is a weak-* closed ideal in a factor and is therefore zero. Thus \(p_1q\ne0\).

For \(0<\varepsilon<1\), the spectral projections
\(p_j=\mathbf1_{(1-\varepsilon,1]}(e_j)\) are nonzero. Their product is nonzero, and on its range the positive product \(\prod_j e_j\) is at least \((1-\varepsilon)^n\). Letting \(\varepsilon\downarrow0\) proves the lower bound one; the reverse bound follows from contractivity.
\end{proof}

\begin{lemma}[Pure-state excision]\label{lem:excision}
Let \(E\) be a unital \(C^*\)-algebra and \(\rho\) a pure state. Given a finite set \(\mathcal F\subset E\) and \(\varepsilon>0\), there is a positive contraction \(f\in E\) with \(\|f\|=1\) such that
\[
 \|fxf-\rho(x)f^2\|<\varepsilon\qquad(x\in\mathcal F).
\]
\end{lemma}
\begin{proof}
Let \((\gamma,H,\xi)\) be the irreducible GNS representation of \(\rho\), with \(\|\xi\|=1\), and put
\[
 L=\{z\in E:\gamma(z)\xi=0\}.
\]
We first show that \(\ker\rho=L+L^*\). The inclusion from right to left is immediate. If \(\rho(x)=0\), write \(\eta=\gamma(x)\xi\perp\xi\). If \(\eta=0\), then \(x\in L\). Otherwise, on \(\operatorname{span}\{\xi,\eta\}\) let \(T\) be the rank-one operator with \(T\xi=\eta\) and \(T\eta=0\); thus \(T^*\xi=0\). Apply the selfadjoint form of Kadison transitivity \cite[Theorem~1]{Kadison} separately to the real and imaginary parts of \(T\). This gives \(y\in E\) such that
\[
 \gamma(y)\xi=\eta,\qquad\gamma(y^*)\xi=0.
\]
Then \(x-y\in L\) and \(y\in L^*\), as desired.

The algebra \(L\cap L^*\) is hereditary. Let \((g_\lambda)\) be a positive contractive approximate identity for it. If \(z\in L\), then \(z^*z\in L\cap L^*\), whence
\[
 \|z(1-g_\lambda)\|^2
 =\|(1-g_\lambda)z^*z(1-g_\lambda)\|\longrightarrow0.
\]
The adjoint assertion holds for elements of \(L^*\). Therefore, with \(f_\lambda=1-g_\lambda\), the equality \(\ker\rho=L+L^*\) gives
\[
 \|f_\lambda(x-\rho(x)1)f_\lambda\|\longrightarrow0
 \qquad(x\in E).
\]
Each \(f_\lambda\) is a positive contraction of norm one because \(\gamma(f_\lambda)\xi=\xi\). One sufficiently large \(\lambda\) works for \(\mathcal F\).
\end{proof}

\begin{proof}[Proof of Theorem~\ref{thm:tensor-intro}]
The closed ideal generated by commutators is all of \(D\); otherwise a nonzero unital commutative quotient would supply a character. Decomposing elements into selfadjoint parts and rescaling, choose selfadjoint contractions \(s_i,t_i\) and elements \(v_i,w_i\), for \(1\leq i\leq m\), such that
\begin{equation}\label{eq:commutator}
 \left\|1-\sum_{i=1}^m v_i[s_i,t_i]w_i\right\|<\tfrac12.
\end{equation}
Put \(s_0=1\), \(k=m+1\), and
\(M=\sum_i\|v_i\|\|w_i\|>0\).
Under any nonzero unital representation \(\sigma\), some represented commutator has norm greater than \(\delta=(2M)^{-1}\). Since
\[
 \|[S,T]\|\leq2\|S-c1\|\|T\|
 \leq\operatorname{diam}(\spc S)
\]
when \(S=S^*\), \(\|T\|\leq1\), and \(c\) is the midpoint of \(\spc S\), the spectrum of some \(\sigma(s_i)\) has diameter at least \(\delta\).

There are pure states \(\rho^0,\rho^1\) of \(\sigma(D)\) taking that element to its two spectral endpoints. Indeed, a state attaining an endpoint extends from the unital abelian algebra \(C^*(1,\sigma(s_i))\), and the states attaining that endpoint form a nonempty weak-* compact face of the state space. An extreme point of this face is pure. The associated real expectation vectors
\[
 v^\alpha=\big(\rho^\alpha(\sigma(s_b))\big)_{b=0}^m
 \in\{1\}\times[-1,1]^m\qquad(\alpha=0,1)
\]
differ by at least \(\delta\) in one coordinate. The set of pairs in \(\big(\{1\}\times[-1,1]^m\big)^2\) with \(\|v^0-v^1\|_\infty\geq\delta\) is compact. Consequently there is \(0<r<1\), independent of \(\sigma\), such that
\begin{equation}\label{eq:angle}
 \left|\left\langle\frac{v^0}{\|v^0\|},\frac{v^1}{\|v^1\|}\right\rangle\right|\leq r.
\end{equation}
To justify the strict bound, two vectors in this compact set cannot be negatively collinear because both first coordinates are one; positive collinearity forces them to coincide, contrary to their separation.

Choose \(G,n\) from Lemma~\ref{lem:gaussian} for \(k,r\) and form the selfadjoint element
\begin{equation}\label{eq:h}
 h=\sum_{b_1,\ldots,b_n=0}^m G_{b_1\ldots b_n}
 s_{b_1}\otimes\cdots\otimes s_{b_n}
 \in D^{\otimes_{\max}n}.
\end{equation}
Let \(\pi\) be any nonzero irreducible representation of this tensor product, with factor restrictions \(\pi_j:D\to\mathcal B(H)\). For each \(j\), choose the two pure states of \(\pi_j(D)\) described above. By \eqref{eq:angle} and Lemma~\ref{lem:gaussian}, choosing one state \(\rho_j\) from each pair makes
\begin{equation}\label{eq:p}
 p=\sum_{b_1,\ldots,b_n=0}^mG_{b_1\ldots b_n}
 \prod_{j=1}^n\rho_j(\pi_j(s_{b_j}))
\end{equation}
strictly positive for one choice and strictly negative for another. Removing the normalizations multiplies each tensor evaluation by a positive number and so preserves its sign.

The von Neumann algebras \(M_j=\pi_j(D)''\) are mutually commuting factors: a central element of any \(M_j\) commutes with every factor image and hence with the irreducible joint representation, so is scalar. For one of the choices in \eqref{eq:p}, use Lemma~\ref{lem:excision} to choose positive norm-one elements \(e_j\in\pi_j(D)\) with
\[
 \|e_j\pi_j(s_b)e_j-\rho_j(\pi_j(s_b))e_j^2\|<\varepsilon
 \quad(0\leq b\leq m).
\]
The product \(e=\prod_j e_j\) is positive and has norm one by Lemma~\ref{lem:factors}. Commuting distinct factors and telescoping their products in \eqref{eq:h} yields
\begin{equation}\label{eq:compression}
 \|e\pi(h)e-pe^2\|
 \leq n\varepsilon\sum_{b_1,\ldots,b_n}|G_{b_1\ldots b_n}|.
\end{equation}
All factors in this estimate are contractions.

If \(p>0\) and \(\pi(h)\leq0\), then \(pe^2-e\pi(h)e\geq pe^2\), so the left side of \eqref{eq:compression} is at least \(p\). Choosing \(\varepsilon\) sufficiently small is a contradiction. The negative choice of \(p\), with its own excisers, similarly rules out \(\pi(h)\geq0\). Thus
\[
 \pi(h_+)\ne0,\qquad\pi((-h)_+)\ne0
\]
for every nonzero irreducible \(\pi\). Each positive part is full: a proper closed ideal containing it would yield a nonzero quotient and hence an irreducible representation killing it. The two positive parts are orthogonal by functional calculus.
\end{proof}

The canonical unital map from the finite tensor power into the infinite maximal tensor power preserves fullness and orthogonality. Thus Theorem~\ref{thm:tensor-intro} also gives the affirmative answer to \cite[Question 6.2]{KR}.
